% ===================================================================
\section{Preguntas generales para la defensa}

Estas preguntas no dependen de una diapositiva concreta.

\begin{preguntas}
\pr{¿Cuál es la contribución original del proyecto?}{Mostrar, con un experimento
controlado, que la sincronización de los ciclos de pago, estimada solo con
transacciones, contiene información sobre el impago que no está en el score clásico ni
en la estructura estática de la red. Incluye el método para estimarla (fase por eventos y
coherencia local) y un control nulo que descarta artefactos.}
\pr{¿Por qué un físico trabaja en riesgo de crédito?}{Porque las herramientas son de la
física: sistemas dinámicos, análisis espectral, osciladores acoplados, transiciones de
fase y avalanchas. La red de pagos es un sistema complejo y estas herramientas permiten
responder cuándo amplifica y cómo se sincroniza.}
\pr{¿El resultado es causal o solo predictivo?}{En el mundo simulado la causa se conoce
porque la pusimos: el desfase genera déficit. El experimento de $\gamma$ muestra que la
señal depende de esa causa. Con datos reales, el método sería predictivo; establecer
causalidad requeriría otro diseño.}
\pr{¿Cómo se evitó la fuga de información?}{De cuatro formas. Las variables solo usan
transacciones anteriores al fin de la observación, y el código rechaza cualquier pago
posterior. Las columnas de verdad del modelo están vetadas por esquema. La validación
separa carteras completas. Y el control nulo muestra mejora cero.}
\pr{¿Por qué no se usó una GCN en la tarea 6?}{Las variables de red ya resumen la
información relevante de los vecinos de cada cliente, y la pregunta es qué tipo de
información aporta. Comparar niveles de variables con el mismo modelo responde esa
pregunta de forma más limpia que un modelo de grafo de caja negra.}
\pr{¿Qué pasa si los parámetros son otros?}{El tamaño de la mejora cambia: crece con
$\gamma$ y se anula sin él. Lo que se mantiene es el patrón: sin mecanismo no hay
mejora, con mecanismo sí. Los parámetros están en \texttt{configs/credit.yaml} y el
experimento completo se vuelve a correr en unos 20 segundos.}
\pr{¿Funciona si el banco no ve toda la red?}{No se probó. Un banco ve solo a sus
clientes y los pagos que pasan por él. Es parte del trabajo futuro: el operador de un
sistema interoperable como Bre-B tendría la vista completa.}
\pr{¿Cuánto tarda todo?}{Las pruebas rápidas, unos 10 segundos. El experimento de la
tarea 6, unos 20 segundos. El conjunto de datos de la tarea 5, unos 20 minutos en ocho
núcleos.}
\end{preguntas}

% ===================================================================
\section{Glosario de siglas, modelos y métodos}

\begin{longtable}{p{0.2\textwidth}p{0.74\textwidth}}
\toprule
\textbf{Término} & \textbf{Significado} \\
\midrule
\endhead
AUC & Área bajo la curva ROC. Probabilidad de que un positivo al azar reciba más riesgo
que un negativo al azar. \\
BA & Barabási–Albert: red con conexión preferencial y nodos concentradores. \\
Bre-B & Sistema de pagos inmediatos del Banco de la República de Colombia. \\
Brier & Error cuadrático medio de una probabilidad predicha. \\
DANE & Departamento Administrativo Nacional de Estadística. \\
DebtRank & Medida de propagación de estrés en redes financieras (Battiston et al.,
2012). \\
ECE & Error de calibración esperado. \\
Eisenberg–Noe & Modelo de liquidación de obligaciones en una red (2001). \\
ER & Erdős–Rényi: red aleatoria con probabilidad de conexión uniforme. \\
Euler–Maruyama & Integrador de ecuaciones diferenciales con ruido. \\
Fiedler & Vector propio de la conectividad algebraica $\lambda_2$; separa comunidades. \\
GCN & Red neuronal convolucional sobre grafos. \\
Gini & $2\,\mathrm{AUC}-1$. Métrica de discriminación usada en banca. \\
GroupKFold & Validación cruzada que mantiene cada grupo en una sola partición. \\
HHI & Índice de Herfindahl–Hirschman: concentración, suma de participaciones al
cuadrado. \\
IC 95~\% & Intervalo de confianza del 95~\%. \\
Kuramoto & Modelo de osciladores de fase acoplados. \\
KS & Estadístico de Kolmogorov–Smirnov: máxima separación entre distribuciones
acumuladas. \\
Laplaciano & Matriz $D-W$ que describe flujos que se conservan en una red. \\
Lorentz & Distribución de Cauchy; da la solución exacta de Kuramoto en la red
completa. \\
NP & Núcleo-periferia. \\
Parquet & Formato de archivo para tablas, comprimido y con tipos fijos. \\
Pareto & Distribución de cola pesada: pocos valores muy grandes. \\
PD & Probabilidad de incumplimiento a un horizonte (aquí 12 meses). \\
Perron–Frobenius & Teorema sobre el valor propio dominante de matrices no negativas. \\
PINN & Red neuronal informada por la física. \\
PR-AUC & Área bajo la curva precisión-exhaustividad. \\
PageRank & Centralidad por recorridos aleatorios sobre la red. \\
$R$ & Índice de amplificación $\beta\rho(W)/\alpha$. \\
$r$, $r_i$ & Parámetro de orden global y coherencia local. \\
RK4 & Runge–Kutta de cuarto orden. \\
ROC & Curva de tasa de verdaderos positivos contra falsos positivos. \\
SBM & Modelo de bloques estocástico: la probabilidad de conexión depende del grupo de
cada nodo. Se usa en las transferencias. \\
SeedSequence & Generador de semillas independientes de NumPy. \\
SHAP & Explicación aditiva de predicciones basada en valores de Shapley. \\
SMMLV & Salario mínimo mensual legal vigente. \\
Spearman & Correlación entre rangos. \\
Watts & Modelo de cascadas por umbral fraccional (2002). \\
XGBoost & Árboles de decisión con \emph{gradient boosting}. \\
YAML & Formato de texto para archivos de configuración. \\
\bottomrule
\end{longtable}
